% Research Report 2 — A direct attempt at the core
% Build: latexmk report-02.tex
\documentclass[11pt,a4paper]{article}

\usepackage{common}

\fancyhf{}
\fancyhead[L]{\small\sffamily Yang--Mills programme \textperiodcentered{} Research Report 2}
\fancyhead[R]{\small\sffamily 2026-09-25}
\fancyfoot[C]{\small\sffamily\thepage}
% Ledger identifiers refer to spec.pdf; they are not links in this document.
\newcommand{\idref}[1]{\textsf{#1}}
\newcommand{\yes}{\textcolor{Green}{\ding{51}}}
\newcommand{\no}{\textcolor{BrickRed}{\ding{55}}}
\newcommand{\open}{\textcolor{Orange}{?}}
\usepackage{pifont}

\hypersetup{pdftitle={Research Report 2: A direct attempt at the core},pdfauthor={Christian Bucher}}

\begin{document}

\begin{center}
{\sffamily\small YANG--MILLS EXISTENCE AND MASS GAP \textperiodcentered{} RESEARCH PROGRAMME}\\[8pt]
{\LARGE\bfseries Research Report 2}\\[4pt]
{\Large A direct attempt at the core:\\candidate mechanisms and the multiscale Bakry--Émery route}\\[10pt]
{\large Christian Bucher}\footnote{Research and drafting assistance: Claude Code (Anthropic). Literature statements were checked against the primary sources. The AI tool is not an author~\cite{COPE23}.}\\[4pt]
{\small 2026-09-25 \textperiodcentered{} companion to \texttt{spec.pdf} v1.2 and Research Report 1}
\end{center}

\begin{keybox}[title=Result in brief]
\textbf{The attempt did not produce a solution. It contains no proof of the mass gap.} It attacks the open core \idref{A-031} directly:
\begin{enumerate}[leftmargin=1.4em]
\item \textbf{Nine candidate mechanisms} are passed through the mandatory sanity tests of the specification (Section~\ref{sec:funnel}). None survives as a complete route. Each is either heuristic, or tied to one dimension, one regime or finite volume, or it reduces to the same open statement.
\item \textbf{The strongest current tool}, the multiscale Bakry--Émery criterion of Bauerschmidt, Bodineau and Dagallier~\cite{BBD24}, is worked through for lattice Yang--Mills (Section~\ref{sec:mbe}). It meets four obstacles. The decisive one is that the criterion needs the renormalized potential to become uniformly convex on large scales, which is mass generation restated. We call this statement \idref{A-043}.
\item \textbf{A precise gauge dilemma} appears along the way: the gauges available on the lattice are either global but non-local, or local but not global. A Polchinski-flow approach needs both properties (\idref{A-042}).
\end{enumerate}
\textbf{Status.} The closure count is unchanged (1 of 9 transitions). The ledger gains five entries, \idref{A-041}--\idref{A-045}.
\end{keybox}

% ======================================================================
\section{Target and filter}\label{sec:target}

The target is \idref{A-031}: there are $c>0$, $a_0$ and $L_0$ such that $\Delta_{a,L}\ge c\,\Lambda$ for all $0<a<a_0$ and $L>L_0$, in physical units, for every compact simple $\mathfrak g$. Research Report 1 showed that this is a bridge across a coupling window. Every candidate mechanism is first passed through the mandatory sanity tests of the specification (\texttt{spec.pdf}, \S12.3), extended by two further columns:
\begin{description}[font=\normalfont\sffamily,leftmargin=1.3cm,labelwidth=1.1cm,align=left]
\item[Ab] it must \emph{fail} for $G=U(1)$, which is free and gapless;
\item[Free] it must fail at $g=0$;
\item[Vol] it must be uniform in the volume;
\item[UV] it must be compatible with the continuum limit $a\to0$ and asymptotic freedom;
\item[All $\mathfrak g$] it must work for every compact simple $\mathfrak g$;
\item[Rig.] it must be rigorous, or plausibly made so.
\end{description}

% ======================================================================
\section{Candidate mechanisms}\label{sec:funnel}

\begin{center}\small
\begin{longtable}{@{}P{4.1cm}cccccc P{4.9cm}@{}}
\toprule
Mechanism & Ab & Free & Vol & UV & All $\mathfrak g$ & Rig. & Verdict\\
\midrule\endhead
\bottomrule\endfoot
C1 Quartic potential, flat directions~\cite{Fey81,Sim83} & \yes & \yes & \no & \open & \yes & \yes & Discrete spectrum only in finite volume; the level spacing is not uniform as $L\to\infty$.\\
C2 Curvature of the orbit space~\cite{Sin81,BV81} & \yes & \yes & \open & \open & \yes & \no & The Ricci curvature needs regularization; no uniform lower bound is known. Report~1, Prop.~3.2 is the lattice analogue and fails at weak coupling.\\
C3 Gauge-invariant variables, Jacobian-induced mass~\cite{KN96,KKN98} & \yes & \yes & \open & \open & \yes & \no & Heuristic and in $2+1$ dimensions; no analogue in $3+1$ dimensions.\\
C4 Strong-coupling expansions~\cite{OSe78,SZZ23} & \yes & \yes & \yes & \no & \yes & \yes & Rigorous, but in the regime opposite to the continuum limit (Report~1).\\
C5 Multiscale Bakry--Émery~\cite{BBD24} & \yes & \yes & \open & \open & \open & \yes & The strongest current tool; see Section~\ref{sec:mbe}. It reduces to \idref{A-043}.\\
C6 RG plus finite-volume criterion plus certification & \yes & \yes & \open & \open & \open & \yes & Report~1: blocked by the coupling window and the dimensional barrier.\\
C7 Large $N$, string duality~\cite{tH74} & --- & --- & \open & \open & \no & \no & Heuristic, and needs $N\to\infty$, whereas the problem fixes $G$.\\
C8 Correlation inequalities & --- & --- & --- & --- & \no & \yes & They rely on properties of scalar or abelian models~\cite[\S6.1]{JW}.\\
C9 Exact solution & --- & --- & --- & --- & --- & \no & \enquote{Appears unlikely} in four dimensions~\cite[\S6.1]{JW}.\\
\end{longtable}
\end{center}
{\small \yes{} passes the test, \no{} fails it, \open{} open, --- not applicable. \enquote{Ab: \yes} means the mechanism correctly gives no gap for $U(1)$.}

\textbf{Reading.} C4 is the only complete rigorous mechanism, and it lives in the wrong regime. C5 and C6 are rigorous frameworks with open inputs. Everything else is heuristic or confined to finite volume, a lower dimension or large $N$. C5 is the one tool that has provably crossed windows of a comparable kind in other models, so we examine it in detail.

% ======================================================================
\section{The multiscale Bakry--Émery route}\label{sec:mbe}

\subsection{The criterion}
Let $\nu_0$ be a probability measure on a finite-dimensional vector space $X$ of the form $\nu_0(d\varphi)\propto e^{-V_0(\varphi)}P_C(d\varphi)$, where $P_C$ is the Gaussian measure with covariance $C=\int_0^\infty\dot C_t\,dt$. Let $V_t$ be the renormalized potential defined by the Polchinski flow~\cite{Pol84}.

\begin{importedtheorem}[{\cite[Thm.~3.6]{BBD24}}; \idref{A-041}]\label{thm:bbd}
Suppose there are real numbers $\dot\lambda_t$, which may be negative, such that
\[
\dot C_t\,\mathrm{Hess}\,V_t(\varphi)\,\dot C_t-\tfrac12\ddot C_t\;\ge\;\dot\lambda_t\,\dot C_t\qquad\text{for all }\varphi\in X,\ t>0,
\]
and put $\lambda_t=\int_0^t\dot\lambda_s\,ds$ and $\frac1\gamma=\int_0^\infty e^{-2\lambda_t}\,dt$. Then $\nu_0$ satisfies the log-Sobolev inequality $\mathrm{Ent}_{\nu_0}[F]\le\frac2\gamma\,\mathbb E_{\nu_0}\bigl[(\nabla\sqrt F)^2_{\dot C_0}\bigr]$.
\end{importedtheorem}

The point of the criterion is that $V_0$ need not be convex: only an \emph{integrated} bound along the flow is required. Temporary non-convexity on finitely many scales costs a finite factor. With this criterion, log-Sobolev inequalities were proved for the continuum sine-Gordon model~\cite{BB20}, for near-critical Ising models~\cite{BD23a} and for $\phi^4_2$ and $\phi^4_3$~\cite{BD23b}. These are models in which pointwise convexity fails.

\subsection{Applying it to lattice Yang--Mills: four obstacles}

\paragraph{O1: compact target and large fields.} The configuration space $G^E$ is a compact manifold, not a vector space with a Gaussian reference. Linear coordinates exist only in charts. For $SU(2)$, $U=\exp(i\,a\cdot\sigma)$ with $|a|<\pi$, and the Haar density is $(\sin|a|/|a|)^2$. In such coordinates $V_0$ is $+\infty$ outside the chart and non-convex near its boundary. This is the lattice form of Ba{\l}aban's large-field problem~\cite{Bal89a,Bal89b}. The Polchinski flow smooths $V_0$, but there is no a priori control of $\mathrm{Hess}\,V_t$ near large fields.

\paragraph{O2: the gauge dilemma (\idref{A-042}).} A Gaussian reference needs a gauge fixing. On a finite lattice, fixing the links of a maximal tree to $\mathbb 1$ is a \emph{global} gauge fixing: the remaining links keep the Wilson weight and the Haar measure. It is not \emph{local}, however. In the Gaussian (weak-coupling) approximation, the temporal-gauge field is a sum of independent electric increments along the time direction. Its variance therefore grows linearly with the time extent, and the reference covariance is not bounded uniformly in the volume. Local covariant gauges avoid this, but by the Gribov ambiguity and Singer's theorem~\cite{Gri78,Sin78} they are not globally defined once large fields matter. The Polchinski framework needs a reference that is both global and local, uniformly in the volume. Ba{\l}aban resolves the analogous problem for his RG with gauges that vary from block to block and from scale to scale~\cite{Bal87}. A gauge-covariant version of the Polchinski flow along these lines does not yet exist.

\paragraph{O3: mass generation as large-scale convexification (\idref{A-043}).} This is the decisive obstacle. For a massless Gaussian reference, $V_0=0$ gives $\lambda_t$ bounded, and $1/\gamma$ grows with the volume: there is no uniform log-Sobolev constant, as there must not be for a gapless theory. For the criterion to give a constant uniform in the volume, the renormalized potential must provide the missing convexity on large scales. That is, $\dot C_t\,\mathrm{Hess}\,V_t\,\dot C_t$ must eventually dominate, uniformly in $\varphi$ and the volume. In a theory whose classical action is scale invariant and whose Gaussian part is massless, this convexity \emph{is} the dynamically generated mass. The mass-gap problem is thereby restated precisely as a convexification statement for the renormalized Yang--Mills potential beyond the crossover scale. It is not solved. Asymptotic freedom controls $V_t$ only while the effective coupling is small, and the required convexity is needed exactly where it is not (Report~1, Obstruction~1).

\paragraph{O4: from a dynamical gap to the physical gap (\idref{A-044}).} A log-Sobolev or Poincaré inequality is a statement about a stochastic dynamics. The mass gap concerns the decay of equilibrium correlations. The generic conversion, via finite speed of propagation, yields a decay rate proportional to the dynamical gap. For the massive Gaussian free field with the standard Langevin dynamics the dynamical gap is $m^2$ while the correlations decay at rate $m$. In lattice units near the continuum limit $m\to0$, so the generic conversion loses the square root and with it the physical scaling. A usable route therefore needs a conversion adapted to the preconditioned Dirichlet form $(\nabla\sqrt F)^2_{\dot C_0}$ of Theorem~\ref{thm:bbd}, or an operator-valued (Brascamp--Lieb-type) bound, for gauge-invariant observables. We found no such statement in the literature for gauge theories.

\subsection{What the route would need}
\begin{proposal}[Multiscale route; recorded as \idref{A-045}, \statusbadge{OPEN}]
The following three statements would imply \idref{A-031}:
\begin{enumerate}[label=(M\arabic*)]
\item a gauge-covariant Polchinski-type flow for lattice Yang--Mills that resolves O1 and O2 uniformly in $a$ and the volume (\idref{A-042});
\item large-scale convexification: along this flow, the criterion of Theorem~\ref{thm:bbd} holds with a log-Sobolev constant of order $(a\Lambda)^{-2}$, uniformly in $a$ and the volume (\idref{A-043});
\item a conversion of that inequality into exponential decay of gauge-invariant correlations at rate $\gtrsim a\Lambda$ in lattice units (\idref{A-044}).
\end{enumerate}
\end{proposal}
The implication is itself only sketched. (M2) is the mass-gap problem in multiscale form; (M1) and (M3) are substantial open problems of their own.

% ======================================================================
\section{Verdict}

A direct attempt with the strongest available tools ends at the same place as Report~1, now stated more sharply. Every rigorous route we can identify needs a non-perturbative statement at and beyond the crossover scale. In the multiscale Bakry--Émery language it is the convexification statement (M2), which is mass generation itself. We did not find a new idea that proves it, and we do not claim one. Any future claim must pass all columns of the filter in Section~\ref{sec:funnel}.

% ======================================================================
\appendix
\section{Deutsche Zusammenfassung}
Der direkte Angriff auf den Kern des Problems hat \textbf{keine Lösung} ergeben. Neun Mechanismen wurden den Pflichttests der Spezifikation unterzogen. Keiner trägt vollständig: Sie sind heuristisch, auf endliches Volumen, eine niedrigere Dimension oder große $N$ beschränkt, liegen im falschen Kopplungsbereich oder führen auf dieselbe offene Aussage zurück. Das stärkste heutige Werkzeug, das mehrskalige Bakry--Émery-Kriterium, hat in anderen Modellen vergleichbare Fenster überwunden. Für Yang--Mills stößt es auf vier Hindernisse. Das entscheidende: Das Kriterium verlangt, dass das renormierte Potential auf großen Skalen gleichmäßig konvex wird. Genau das ist die Massenerzeugung, also das Mass-Gap-Problem in anderer Form. Unterwegs zeigt sich ein präzises Eichdilemma: Die verfügbaren Eichungen sind entweder global, aber nicht lokal (axiale Eichung), oder lokal, aber nicht global (Gribov/Singer). Eine neue Idee, die das überwindet, haben wir nicht gefunden.

\begin{thebibliography}{KKN98}
\small
\bibitem[JW]{JW} A.~Jaffe, E.~Witten, \emph{Quantum Yang--Mills theory}, official problem description, Clay Mathematics Institute (2000).
\bibitem[BBD24]{BBD24} R.~Bauerschmidt, T.~Bodineau, B.~Dagallier, Stochastic dynamics and the Polchinski equation: an introduction, \emph{Probab. Surveys} 21 (2024). doi:10.1214/24-PS27, arXiv:2307.07619.
\bibitem[BB20]{BB20} R.~Bauerschmidt, T.~Bodineau, Log-Sobolev inequality for the continuum sine-Gordon model, \emph{Comm. Pure Appl. Math.} 74 (2021) 2064--2113. doi:10.1002/cpa.21926
\bibitem[BD23a]{BD23a} R.~Bauerschmidt, B.~Dagallier, Log-Sobolev inequality for near critical Ising models, \emph{Comm. Pure Appl. Math.} 77 (2024) 2568--2576. doi:10.1002/cpa.22172
\bibitem[BD23b]{BD23b} R.~Bauerschmidt, B.~Dagallier, Log-Sobolev inequality for the $\varphi^4_2$ and $\varphi^4_3$ measures, \emph{Comm. Pure Appl. Math.} 77 (2024) 2579--2612. doi:10.1002/cpa.22173
\bibitem[Pol84]{Pol84} J.~Polchinski, Renormalization and effective lagrangians, \emph{Nucl. Phys. B} 231 (1984) 269--295. doi:10.1016/0550-3213(84)90287-6
\bibitem[Fey81]{Fey81} R.~P.~Feynman, The qualitative behavior of Yang--Mills theory in 2+1 dimensions, \emph{Nucl. Phys. B} 188 (1981) 479--512. doi:10.1016/0550-3213(81)90005-5
\bibitem[Sim83]{Sim83} B.~Simon, Some quantum operators with discrete spectrum but classically continuous spectrum, \emph{Ann. Phys.} 146 (1983) 209--220. doi:10.1016/0003-4916(83)90057-X
\bibitem[Sin81]{Sin81} I.~M.~Singer, The geometry of the orbit space for non-abelian gauge theories, \emph{Phys. Scripta} 24 (1981) 817--820. doi:10.1088/0031-8949/24/5/002
\bibitem[BV81]{BV81} O.~Babelon, C.~M.~Viallet, The Riemannian geometry of the configuration space of gauge theories, \emph{Commun. Math. Phys.} 81 (1981) 515--525. doi:10.1007/BF01208272
\bibitem[KN96]{KN96} D.~Karabali, V.~P.~Nair, A gauge-invariant Hamiltonian analysis for non-abelian gauge theories in (2+1) dimensions, \emph{Nucl. Phys. B} 464 (1996) 135--152. doi:10.1016/0550-3213(96)00034-X
\bibitem[KKN98]{KKN98} D.~Karabali, C.~Kim, V.~P.~Nair, Planar Yang--Mills theory: Hamiltonian, regulators and mass gap, \emph{Nucl. Phys. B} 524 (1998) 661--694. doi:10.1016/S0550-3213(98)00309-5
\bibitem[OSe78]{OSe78} K.~Osterwalder, E.~Seiler, Gauge field theories on a lattice, \emph{Ann. Phys.} 110 (1978) 440--471. doi:10.1016/0003-4916(78)90039-8
\bibitem[SZZ23]{SZZ23} H.~Shen, R.~Zhu, X.~Zhu, A stochastic analysis approach to lattice Yang--Mills at strong coupling, \emph{Commun. Math. Phys.} 400 (2023) 805--851. doi:10.1007/s00220-022-04609-1
\bibitem[tH74]{tH74} G.~'t~Hooft, A planar diagram theory for strong interactions, \emph{Nucl. Phys. B} 72 (1974) 461--473.
\bibitem[Gri78]{Gri78} V.~N.~Gribov, Quantization of non-abelian gauge theories, \emph{Nucl. Phys. B} 139 (1978) 1--19. doi:10.1016/0550-3213(78)90175-X
\bibitem[Sin78]{Sin78} I.~M.~Singer, Some remarks on the Gribov ambiguity, \emph{Commun. Math. Phys.} 60 (1978) 7--12. doi:10.1007/BF01609471
\bibitem[Bal87]{Bal87} T.~Ba{\l}aban, Renormalization group approach to lattice gauge field theories I, \emph{Commun. Math. Phys.} 109 (1987) 249--301.
\bibitem[Bal89a]{Bal89a} T.~Ba{\l}aban, Large field renormalization I, \emph{Commun. Math. Phys.} 122 (1989) 175--202. doi:10.1007/BF01257412
\bibitem[Bal89b]{Bal89b} T.~Ba{\l}aban, Large field renormalization II, \emph{Commun. Math. Phys.} 122 (1989) 355--392. doi:10.1007/BF01238433
\bibitem[COPE23]{COPE23} Committee on Publication Ethics, \emph{Authorship and AI tools}, COPE position statement (2023).
\end{thebibliography}

\end{document}
